% Original Prism Journal of Medicinal Chemistry Perspective draft (MIT).
% This is an interpretive writing skeleton, not a final journal layout.
\documentclass[letterpaper,12pt]{article}
\usepackage[margin=1in]{geometry}
\usepackage{amsmath,amssymb,graphicx,array,booktabs,tabularx,siunitx,newfloat}
\usepackage[version=4]{mhchem}
\usepackage{setspace}
% Select package-mode settings explicitly rather than assuming a journal cfg.
\usepackage[articletitle=true,doi=true,super=true,biblabel=brackets,
  maxauthors=15,etalmode=firstonly]{achemso}
\usepackage[hidelinks]{hyperref}
\providecommand{\doi}[1]{\href{https://doi.org/#1}{\nolinkurl{#1}}}
\DeclareFloatingEnvironment[fileext=los,name=Scheme,placement=htbp]{scheme}
\DeclareSIUnit{\molar}{M}
\sisetup{uncertainty-mode=separate}
\setcounter{secnumdepth}{0}
\setlength{\emergencystretch}{2em}
\doublespacing
\clubpenalty=10000
\widowpenalty=10000
\displaywidowpenalty=10000
\raggedbottom

\title{Perspective manuscript title}
\author{First Author\textsuperscript{*} and Second Author\\
  \small Department, Institution, City, Country\\
  \small \textsuperscript{*}Correspondence: \texttt{first.author@example.org}}
\date{}

\begin{document}
\maketitle
\begin{abstract}
Replace this paragraph with a self-contained English abstract of
approximately 150 words. Introduce the medicinal chemistry subject, state
the central interpretation and explain the educational value of the
Perspective. Summarize only the themes and conclusions developed from
verified evidence in the manuscript. Approximately 150 words is a drafting
prompt; confirm the current guidance rather than treating this template as
an automatic compliance check. This starter contains synthetic displays
and writing instructions only. It does not document a literature search,
experimental study or author achievement. A Perspective offers critical
insight and teaching moments rather than claiming exhaustive coverage.
Replace all placeholders before submission and discuss an unsolicited topic
with the journal's editor.
\end{abstract}

\section{Introduction}
Define the question, intended readership and boundaries of the Perspective.
Explain why the topic matters and identify the central insight or position
to be developed. Perspectives are usually invited; prospective authors
should contact the editor to discuss suitability before submitting an
unsolicited manuscript. The references here demonstrate citation handling
only, not evidence for the synthetic displays\cite{peng2011reproducible,wilkinson2016fair}.

\section{Conceptual framework and teaching moments}
\subsection{Explain the organizing principle}
Introduce the concepts needed to follow the argument. Develop concrete
teaching moments from verified examples, explaining what the reader should
learn and where a principle ceases to apply. Chemical notation such as
\ce{CH3}, \ce{OH} and \ce{C=O} is demonstrated here without assigning it to
a synthesized, tested or literature-reported compound.

Scheme~\ref{sch:concept} is a synthetic conceptual slot. Replace it with
original artwork that supports the Perspective's actual argument, with
appropriate source attribution where needed.
\begin{scheme}[htbp]
\centering
\fbox{\parbox[c][21mm][c]{0.76\linewidth}{\centering
\ce{A -> B}\\[4pt]
Conceptual scheme slot: original artwork}}
\caption{Synthetic scheme for formatting only. A and B are placeholder
labels, not identified molecules or an established transformation. Supply
an accurate explanation when replacing the illustration.}
\label{sch:concept}
\end{scheme}

\subsection{Distinguish observation and interpretation}
Identify the evidence for each teaching point. Separate reported
observations, computational predictions, the authors' interpretation and
open questions. Do not describe a comparison or screening procedure unless
the authors actually performed it and can explain its scope.

\section{Evidence, comparisons and limitations}
\subsection{Compare like with like}
Table~\ref{tab:comparison} demonstrates an SAR-style comparison layout.
Its rows were invented for typesetting and are not extracted from papers.
Real comparisons must identify their sources and account for assay
conditions, units, endpoints, uncertainty and missing information.

\begin{table}[htbp]
\centering\small
\caption{Synthetic comparative values, not a literature-derived SAR.}
\label{tab:comparison}
\begin{tabularx}{\linewidth}{@{}>{\raggedright\arraybackslash}X
 c S[table-format=2.1(2)] S[table-format=1.0]@{}}
\toprule
Placeholder & {Example group} & {$\mathrm{IC}_{50}$ / \unit{\nano\molar}} & {$n$}\\
\midrule
\textbf{1} & \ce{H}   & 12.4 +- 1.2 & 3\\
\textbf{2} & \ce{CH3} & 28.6 +- 2.4 & 3\\
\textbf{3} & \ce{F}   & \multicolumn{1}{c}{$>1000$} & 3\\
\bottomrule
\end{tabularx}
\par\smallskip\footnotesize\raggedright
\textsuperscript{a}All entries, compound labels and counts are invented.
The uncertainties illustrate mean $\pm$ SD only. Cite the actual evidence,
define its conditions and uncertainty, and preserve censored or unavailable
values without treating them as zero.
\end{table}

\subsection{Connect claims to graphics}
Figure~\ref{fig:comparison} provides two related panel slots. An effective
Perspective graphic makes the argument easier to understand while keeping
the boundaries of the evidence visible.
\begin{figure}[htbp]
\centering
\begin{minipage}[t]{0.47\linewidth}
\centering\textbf{a}\par\smallskip
\fbox{\parbox[c][28mm][c]{0.89\linewidth}{\centering
Conceptual graphic slot\\[4pt]Original vector artwork}}
\end{minipage}\hfill
\begin{minipage}[t]{0.47\linewidth}
\centering\textbf{b}\par\smallskip
\fbox{\parbox[c][28mm][c]{0.89\linewidth}{\centering
Evidence graphic slot\\[4pt]Readable final labels}}
\end{minipage}
\caption{Synthetic comparison layout. (a) Conceptual illustration slot.
(b) Evidence illustration slot. Neither panel reports a molecule,
measurement or scientific conclusion. Replace each slot with relevant
original artwork and explain actual symbols, sources, conditions and
limitations. Preserve readable labels and vector quality at final size.
Plan the page using the full caption height, and keep panel letters
consistent with references in the text.}
\label{fig:comparison}
\end{figure}

\subsection{Make quantitative reasoning explicit}
If quantitative expressions help the argument, define their variables,
dimensions and assumptions. Equations~\eqref{eq:pic50} and~\eqref{eq:delta}
demonstrate aligned notation only:
\begin{align}
 p\mathrm{IC}_{50} &= -\log_{10}\!\left(
 \frac{\mathrm{IC}_{50}}{\qty{1}{\mole\per\litre}}\right),\label{eq:pic50}\\
 \Delta p\mathrm{IC}_{50} &= p\mathrm{IC}_{50,\mathrm{A}}
                         -p\mathrm{IC}_{50,\mathrm{B}}.\label{eq:delta}
\end{align}
The concentration ratio is dimensionless. These expressions do not turn
the synthetic table into evidence. Explain relevant comparability limits
before interpreting a difference between published measurements.

\section{Critical discussion and strategic directions}
Develop the authors' reasoned assessment of the evidence. Explain
agreements, contradictions, limitations and competing interpretations.
Identify useful strategic directions and what evidence would be needed to
distinguish them, without reporting proposed work as a completed result.

\section{Conclusions and outlook}
State the principal teaching points and the perspective justified by the
discussion. Keep recommendations proportional to the evidence, acknowledge
the limits of the coverage and identify consequential open questions.

\section*{Supporting Information}
Describe the actual supporting files, if any. The neutral independent
\texttt{supplementary.tex} provides S-numbered display and documentation
slots; it does not assume that experiments or a literature search occurred.
\section*{Data and Code Availability}
Provide an accurate statement appropriate to the actual content and any
data or code supplied by the authors.
\section*{Author Contributions}
Provide the authors' confirmed contributions.
\section*{Acknowledgments}
Provide confirmed funding and acknowledgments, or remove if not applicable.
\section*{Conflict of Interest}
Provide the authors' confirmed statement.

% Keep each author's factual biography immediately before the references.
\section*{Author Biographies}
\noindent\textbf{First Author.} Replace this instruction with a factual
biography of fewer than 125 words, supplied or approved by the author.
Do not invent qualifications, positions, achievements or research history.

\noindent\textbf{Second Author.} Supply an accurate biography of fewer
than 125 words. Repeat the author-provided format for additional authors.

% achemso provides the bibliography style; retain only relevant references.
\bibliography{references}

\clearpage
\section*{Table of Contents graphic}
\noindent For Table of Contents Only
\begin{center}
% The 80 by 42 mm interior leaves room for the frame within the ACS area.
\fbox{\parbox[c][42mm][c]{80mm}{\centering
Original TOC graphic slot\\[6pt]
Replace with the authors' artwork}}
\end{center}
\end{document}
